\chapter{Fundamentacao Teorica}
\label{chap_fundamentacao_teorica}

Um automato celular elementar e definido por uma malha unidimensional binaria e uma regra de transicao local baseada na vizinhanca de tres celulas. Existem 8 padroes possiveis de vizinhanca, e cada padrao produz 0 ou 1 na proxima etapa temporal.

Com isso, existem $2^8=256$ regras possiveis, numeradas de 0 a 255, conhecidas como regras de Wolfram \cite{wolfram1983,wolfram2002}.

As regras podem ser agrupadas em quatro classes qualitativas:
\begin{itemize}
    \item Classe 1 (homogeneo): evolui para estado uniforme.
    \item Classe 2 (periodico): evolui para padroes repetitivos.
    \item Classe 3 (caotico): evolucao aparentemente aleatoria.
    \item Classe 4 (complexo): estruturas localizadas e interacoes ricas.
\end{itemize}

Essa classificacao e amplamente usada para discutir emergencia de complexidade em sistemas discretos \cite{chopard1998,wolfram2002}.
